\scanpage{100}
\uline{Cor} \quad Soit\deleted{\unsure{nt}} $\alpha\in t_0$, $s\in\overline E_\alpha$, alors si $\deleted{\beta,\gamma\in t_0\setminus\{\alpha\},\ \unsure{\beta}\ne\beta,\ \beta\ne\gamma}\deleted{\unsure{\in t_0,\ \beta\ne\alpha}}\beta\in t_0\setminus\{\alpha\}$ l'ens $\widetilde E_\beta(s)$ \inserted{$\subset\overline E_\beta$} des $t\in\widetilde E_\beta$ \deleted{\unsure{$t$}} liés à $s$ est une \uline{section} de $\widetilde E_\beta$ sur $E_\beta=\widetilde E_\beta/\sigma_\beta$. Soit \inserted{de $=$} $\overline{\unsure{a|\gamma}}\deleted{\ne\alpha,\ \ne\beta}\inserted{\in\unsure{\widetilde t|t_0}\setminus\{\alpha,\beta\}}$, $(\gamma\ne\beta)$, d'où $\overline E_\gamma(s)$ section de $\widetilde E_\gamma$ sur $E_\gamma$, la relation entre $t\in\widetilde E_\beta(s)$ et $t'\in\widetilde E_\gamma(s)$ ``$t$ et $t'$ sont liés i.e. $(t,t')\in\widetilde\Gamma_{\beta,\gamma}$'' est une bijection \deleted{entre $\widetilde E_\beta(s)$ et}
\[
u_s^{\gamma,\beta}:\widetilde E_\beta(s)\xrightarrow{\sim}\widetilde E_\gamma(s)
\]

\uline{Cor} \quad Si \deleted{\unsure{on|Si}} $\widetilde E_\alpha\ne\varnothing$, alors $\operatorname{Card}\deleted{\unsure{\sim}}E_\beta=\operatorname{Card}E_\gamma$ donc $\operatorname{Card}\widetilde E_\beta=\operatorname{Card}\widetilde E_\gamma$ ; si ce \deleted{\unsure{dernier}} cardinal est $\ne0$ i.e $\widetilde E_\beta\deleted{\unsure{E_\gamma}}\ne\varnothing$ $(\deleted{\unsure{\longleftarrow|\longleftrightarrow}}\Longleftrightarrow\widetilde E_\gamma\ne\varnothing)$. Alors $\widetilde E_\alpha,\widetilde E_\beta,\widetilde E_\gamma$ ont $=$ cardinal (et de $=$ $E_\alpha,E_\beta,E_\gamma$ ont $=$ cardinal).

\uline{NB} \inserted{\unsure{$\times$}}Évidemt, si $\widetilde E_\alpha=\varnothing\quad\forall\alpha\in t_0$, les axiomes \deleted{\unsure{envisagés}} 1 et 2 sont trivialemt satisfaits. De $\overline{\mathrm m}$ si $\widetilde E_\alpha\ne\varnothing$, et $\widetilde E_\beta,\widetilde E_\gamma=\varnothing$ (donc $\widetilde\Gamma=\varnothing$, donc pas d' \deleted{flèches} arêtes sup transverses). Alors \inserted{le graphe} $S$ est \deleted{\unsure{formé}} de la forme
\[
S=\{\alpha\}\amalg\widetilde S'\deleted{\unsure{a}}\qquad\text{où }\widetilde S'=S-\{\alpha\}
\]
est muni d'une invol. $\sigma_{S'}$ sans pt fixe, et les arêtes sont les $\{\alpha,s\}$ $(s\in S')$ et les $(s,\sigma_{S'}s)$ $(s\in\widetilde S')$, donc
\[
A\simeq\widetilde S'\amalg\underbrace{\widetilde S'/\sigma_{S'}}_{S'}.
\]
\deleted{\unsure{Ma|Mais}}
% archive sheet number: 99
